% !TEX program = xelatex
\documentclass[11pt,a4paper]{article}

\usepackage[UTF8]{ctex}
\usepackage[margin=2.2cm]{geometry}
\usepackage{amsmath, amssymb, amsthm}
\usepackage{listings}
\usepackage{xcolor}
\usepackage{tikz}
\usepackage{booktabs}
\usepackage{multirow}
\usepackage{hyperref}
\usetikzlibrary{positioning, arrows.meta, shapes.geometric, fit, calc, decorations.pathreplacing}

\hypersetup{
  colorlinks=true,
  linkcolor=blue!60!black,
  urlcolor=teal!70!black,
}

\lstdefinelanguage{Rust}{
  morekeywords={fn,let,mut,pub,struct,impl,use,mod,self,Self,Option,Some,None,return,if,else,for,while,loop,match,break,continue,as,where,trait,type,enum,move,ref,const,static,unsafe,extern,crate,super,dyn,async,await,box,in},
  sensitive=true,
  morecomment=[l]{//},
  morecomment=[s]{/*}{*/},
  morestring=[b]",
}

\lstset{
  language=Rust,
  basicstyle=\ttfamily\footnotesize,
  keywordstyle=\color{blue!60!black}\bfseries,
  commentstyle=\color{green!40!black}\itshape,
  stringstyle=\color{red!60!black},
  numbers=left,
  numberstyle=\tiny\color{gray},
  numbersep=8pt,
  frame=single,
  framerule=0.4pt,
  rulecolor=\color{gray!40},
  breaklines=true,
  breakatwhitespace=true,
  showstringspaces=false,
  tabsize=4,
  columns=fullflexible,
  keepspaces=true,
}

\title{\textbf{halfedge\,: 鲁棒几何谓词模块设计文档} \\ \large \texttt{src/predicates.rs}}
\author{halfedge 库}
\date{2026-07-05}

\begin{document}
\maketitle
\tableofcontents

\section{模块定位}

\texttt{predicates.rs} 是 \texttt{halfedge} 库的\textbf{鲁棒几何谓词}模块，
基于 \href{https://docs.rs/robust}{\texttt{robust}} crate
（Jonathan Richard Shewchuk 1997 论文的 Rust 直译实现）提供\textbf{精确的}
几何方向谓词：

\begin{itemize}
  \item \texttt{orient2d}：2D 三点方向判定（CCW / CW / 共线）
  \item \texttt{orient3d}：3D 四点方向判定（在平面下方 / 上方 / 共面）
  \item \texttt{incircle}：2D 四点共圆判定（在圆内 / 圆外 / 共圆）
  \item \texttt{insphere}：3D 五点共球判定（在球内 / 球外 / 共球）
\end{itemize}

以及一批 \texttt{bool} / 标量包装函数（\texttt{is\_ccw2d}、
\texttt{point\_in\_triangle\_2d}、\texttt{tet\_signed\_volume}、
\texttt{is\_triangle\_degenerate\_3d} 等）。

\textbf{使用场景}：布尔运算的射线-三角形求交、三角剖分的凸性 / 内点判定、
拉伸操作的侧面退化检测、网格自检的退化面识别、网格体积计算等。
任何需要\textbf{符号判定}的几何运算都应通过本模块完成。

\begin{center}
\renewcommand{\arraystretch}{1.18}
\begin{tabular}{@{}lll@{}}
  \toprule
  \textbf{类别} & \textbf{API} & \textbf{语义} \\
  \midrule
  \multirow{4}{*}{核心谓词}
    & \texttt{orient2d(a, b, c)}                 & 2D 三点方向（CCW $>$ 0 / CW $<$ 0 / 共线 $=$ 0） \\
    & \texttt{orient3d(a, b, c, d)}              & 3D 四点方向（d 在 abc 平面下方 $>$ 0 / 上方 $<$ 0 / 共面 $=$ 0） \\
    & \texttt{incircle(a, b, c, d)}              & 2D 共圆（d 在 abc 外接圆内 $>$ 0；要求 abc 呈 CCW） \\
    & \texttt{insphere(a, b, c, d, e)}           & 3D 共球（e 在 abcd 外接球内 $>$ 0；要求 abcd 正方向） \\
  \midrule
  \multirow{5}{*}{2D 便利}
    & \texttt{signed\_area2d(a, b, c)}           & $\tfrac{1}{2}\,\text{orient2d}(a,b,c)$（CCW 为正） \\
    & \texttt{triangle\_area2d(a, b, c)}         & 2D 三角形无符号面积 \\
    & \texttt{is\_ccw2d(a, b, c)}                & \texttt{orient2d $>$ 0} \\
    & \texttt{is\_collinear2d(a, b, c)}          & \texttt{orient2d $=$ 0} \\
    & \texttt{is\_convex\_vertex2d(a, b, c)}     & \texttt{orient2d $>$ 0}（ear clipping 凸性判定） \\
  \midrule
  \multirow{2}{*}{2D 点-三角形}
    & \texttt{point\_in\_triangle\_2d(p, a, b, c)} & 4 次 \texttt{orient2d} 同号判定 \\
    & 退化三角形（共线）                          & 返回 \texttt{false} \\
  \midrule
  \multirow{3}{*}{3D 便利}
    & \texttt{tet\_signed\_volume(a, b, c, d)}   & $-\tfrac{1}{6}\,\text{orient3d}(a,b,c,d)$（CCW 朝外为正） \\
    & \texttt{is\_coplanar(a, b, c, d)}          & \texttt{orient3d $=$ 0} \\
    & \texttt{is\_triangle\_degenerate\_3d(a, b, c)} & 主轴投影 + \texttt{orient2d} 精确共线判定 \\
  \bottomrule
\end{tabular}
\end{center}

\section{为什么需要鲁棒谓词}

普通 \texttt{f64} 浮点运算在\textbf{退化情况}（共线、共面、共圆、共球）
附近会因舍入误差给出错误符号，导致依赖符号判定的算法在边界情况下失败：

\begin{itemize}
  \item 布尔运算：点刚好在三角形边上时被误判为内 / 外；
  \item 三角剖分：共线点序列被误判为 CCW / CW，导致 ear clipping 死循环；
  \item 内外判定：射线刚好穿过三角形顶点时方向符号错误；
  \item 退化检测：浮点面积阈值 $10^{-12}$ 对极小但不退化的三角形误报。
\end{itemize}

\textbf{Shewchuk 自适应精度算法}的核心思想是：

\begin{enumerate}
  \item 先用普通 \texttt{f64} 计算快速得到一个近似符号；
  \item 同时计算一个\textbf{误差上界} $\varepsilon$；
  \item 若 $|result| > \varepsilon$，符号确定，直接返回；
  \item 否则启动\textbf{扩展精度}（adaptive precision）重新计算，
        使用 \textbf{expanded double}（将一个 \texttt{f64} 拆为
        非重叠序列的累加形式）消除舍入误差；
  \item 最终保证返回值的\textbf{符号是精确的}（值仍可能不精确，
        但符号决定几何关系，已经足够）。
\end{enumerate}

\textbf{性能}：非退化情况下自适应算法与朴素 \texttt{f64} 实现相当
（仅多一次误差估计）；仅在接近退化时才付出扩展精度的代价。

\section{符号约定}

\subsection{orient2d}

\[
  \text{orient2d}(a, b, c) =
  \det\begin{pmatrix}
    b_x - a_x & b_y - a_y \\
    c_x - a_x & c_y - a_y
  \end{pmatrix}
  = (b_x - a_x)(c_y - a_y) - (b_y - a_y)(c_x - a_x).
\]

\begin{itemize}
  \item $> 0$：$c$ 在有向线 $a \to b$ 的\textbf{左侧}（三点呈 CCW）；
  \item $< 0$：$c$ 在右侧（CW）；
  \item $= 0$：三点共线。
\end{itemize}
返回值绝对值 $=$ 三角形 $abc$ 有向面积的 2 倍。

\subsection{orient3d（Shewchuk 约定）}

\[
  \text{orient3d}(a, b, c, d) =
  \det\begin{pmatrix}
    b_x - a_x & b_y - a_y & b_z - a_z \\
    c_x - a_x & c_y - a_y & c_z - a_z \\
    d_x - a_x & d_y - a_y & d_z - a_z
  \end{pmatrix}.
\]

\textbf{Shewchuk 约定}：从 $abc$ 呈 CCW 顺序的一侧（"上方"）看，
\begin{itemize}
  \item $> 0$：$d$ 在平面 $abc$ 的\textbf{下方}；
  \item $< 0$：$d$ 在\textbf{上方}；
  \item $= 0$：四点共面。
\end{itemize}
返回值绝对值 $=$ 四面体 $abcd$ 有向体积的 6 倍。

\subsection{tet\_signed\_volume（几何约定）}

几何上"\textbf{CCW 朝外为正体积}"的约定与 Shewchuk 相反：

\[
  \text{tet\_signed\_volume}(a, b, c, d) =
  -\frac{1}{6}\,\text{orient3d}(a, b, c, d).
\]

\textbf{取负原因}：当 $abc$ 按右手定则给出朝外法向时，
$d$ 在四面体内部（即 $abc$ 平面"上方"，朝向法向相反一侧），
此时 Shewchuk 的 \texttt{orient3d} 返回负值，
而几何约定希望该四面体体积为正，故取负。

\subsection{incircle / insphere}

\texttt{incircle} 要求 $a, b, c$ 呈 CCW 顺序；\texttt{insphere} 要求
$a, b, c, d$ 呈正方向（即 \texttt{orient3d(a, b, c, d) $>$ 0}）。
若输入违反该前置条件，结果符号反转但仍一致。

\section{核心谓词实现}

\texttt{predicates.rs} 仅做 API 包装，直接委托给 \texttt{robust} crate：

\begin{lstlisting}[language=Rust]
use robust::{Coord, Coord3D};

#[inline]
pub fn orient2d(a: [f64; 2], b: [f64; 2], c: [f64; 2]) -> f64 {
    robust::orient2d(
        Coord { x: a[0], y: a[1] },
        Coord { x: b[0], y: b[1] },
        Coord { x: c[0], y: c[1] },
    )
}

#[inline]
pub fn orient3d(a: [f64; 3], b: [f64; 3], c: [f64; 3], d: [f64; 3]) -> f64 {
    robust::orient3d(
        Coord3D { x: a[0], y: a[1], z: a[2] },
        Coord3D { x: b[0], y: b[1], z: b[2] },
        Coord3D { x: c[0], y: c[1], z: c[2] },
        Coord3D { x: d[0], y: d[1], z: d[2] },
    )
}
\end{lstlisting}

\texttt{robust} crate 是 Shewchuk 1997 论文
\textit{Adaptive Precision Floating-Point Arithmetic and Fast Robust
Geometric Predicates} 的 Rust 直译实现，由 georust 维护，
crates.io 累计下载 490 万+，许可 MIT OR Apache-2.0，
与 \texttt{halfedge} 双许可兼容。

\section{点在三角形内判定}

\texttt{point\_in\_triangle\_2d(p, a, b, c)} 用 4 次 \texttt{orient2d}
同号判定：若 $abc$ 为 CCW，则 $p$ 在内部（含边界）当且仅当 $p$ 相对
三条有向边 $ab$、$bc$、$ca$ 均在左侧或边上（\texttt{orient $\ge 0$}）。

\begin{lstlisting}[language=Rust]
pub fn point_in_triangle_2d(
    p: [f64; 2], a: [f64; 2], b: [f64; 2], c: [f64; 2],
) -> bool {
    let o1 = orient2d(a, b, p);
    let o2 = orient2d(b, c, p);
    let o3 = orient2d(c, a, p);
    let tri_sign = orient2d(a, b, c);
    if tri_sign == 0.0 { return false; }  // 退化三角形
    if tri_sign > 0.0 {  // CCW
        o1 >= 0.0 && o2 >= 0.0 && o3 >= 0.0
    } else {             // CW
        o1 <= 0.0 && o2 <= 0.0 && o3 <= 0.0
    }
}
\end{lstlisting}

\textbf{退化处理}：三角形 $abc$ 共线（\texttt{tri\_sign $=$ 0}）时
返回 \texttt{false}，避免后续 \texttt{>= 0 / <= 0} 永真导致误判。

\section{3D 三角形退化判定}

\texttt{is\_triangle\_degenerate\_3d(a, b, c)} 精确判定 3D 三点是否共线
（含重合）。直接对 3D 三点用 \texttt{orient3d} 无法判共线（只能判共面），
故采用\textbf{主轴投影 + 2D 谓词}的两步法：

\subsection{算法}

\begin{enumerate}
  \item 计算法向 $\vec{n} = (\vec{b} - \vec{a}) \times (\vec{c} - \vec{a})$
        的三个分量绝对值 $|n_x|, |n_y|, |n_z|$；
  \item 选取\textbf{最大分量轴为丢弃轴}，将三角形投影到其余两轴构成的
        2D 平面（保证投影面积最大，避免再次退化）；
  \item 在 2D 平面上调用 \texttt{orient2d} 精确判定共线：
        \texttt{orient2d $=$ 0} $\iff$ 三点共线。
\end{enumerate}

\subsection{投影轴选择公式}

\[
  \text{drop\_axis} =
  \begin{cases}
    x, & |n_x| \ge |n_y| \,\land\, |n_x| \ge |n_z| \\
    y, & |n_y| \ge |n_z| \quad\text{（否则）} \\
    z, & \text{否则}
  \end{cases}
\]

对应投影：
\begin{itemize}
  \item 丢弃 $x$：$(x, y, z) \mapsto (y, z)$
  \item 丢弃 $y$：$(x, y, z) \mapsto (x, z)$
  \item 丢弃 $z$：$(x, y, z) \mapsto (x, y)$
\end{itemize}

\subsection{流程图}

\begin{center}
\resizebox{\textwidth}{!}{%
\begin{tikzpicture}[
  node distance=0.5cm,
  proc/.style={draw, rounded corners=2pt, minimum width=11cm, minimum height=0.7cm, align=center, font=\small},
  dec/.style={diamond, draw, aspect=2.8, fill=orange!12, inner sep=1pt, font=\small, align=center, minimum width=4.5cm},
  op/.style={-Stealth, thick},
  startend/.style={draw, rounded corners=10pt, minimum width=2.8cm, minimum height=0.7cm, font=\small, fill=gray!12},
  term/.style={draw, rounded corners=2pt, fill=green!12, minimum width=4cm, minimum height=0.7cm, font=\small, align=center}
]
  \node[startend] (s) {开始：3D 三点 $a, b, c$};
  \node[proc, below=of s, fill=blue!8] (p1) {计算 $\vec{ab} = b - a$，$\vec{ac} = c - a$，叉积分量 $n_x, n_y, n_z$};
  \node[proc, below=of p1, fill=blue!4] (p2) {取绝对值 $|n_x|, |n_y|, |n_z|$};
  \node[dec, below=of p2] (d1) {$|n_x| \ge |n_y|$ 且 $|n_x| \ge |n_z|$？};
  \node[proc, right=1.0cm of d1, fill=yellow!10] (p3) {丢弃 $x$ 轴，投影到 $yz$ 平面};
  \node[dec, below=1.0cm of d1] (d2) {$|n_y| \ge |n_z|$？};
  \node[proc, below=of d2, fill=yellow!10] (p4) {丢弃 $y$ 轴，投影到 $xz$ 平面};
  \node[proc, right=1.0cm of d2, fill=yellow!10] (p5) {丢弃 $z$ 轴，投影到 $xy$ 平面};
  \node[proc, below=2.0cm of d2, fill=blue!4] (p6) {在 2D 平面调用 \texttt{orient2d}（Shewchuk 自适应精度）};
  \node[dec, below=of p6] (d3) {\texttt{orient2d $=$ 0}？};
  \node[term, right=1.6cm of d3] (t1) {返回 \texttt{true}（退化）};
  \node[term, left=1.6cm of d3] (t2) {返回 \texttt{false}（非退化）};
  \draw[op] (s) -- (p1);
  \draw[op] (p1) -- (p2);
  \draw[op] (p2) -- (d1);
  \draw[op] (d1) -- node[above, font=\scriptsize] {是} (p3);
  \draw[op] (d1) -- node[right, font=\scriptsize] {否} (d2);
  \draw[op] (d2) -- node[right, font=\scriptsize] {是} (p4);
  \draw[op] (d2) -- node[above, font=\scriptsize] {否} (p5);
  \draw[op] (p3.south) |- (p6.east);
  \draw[op] (p4) -- (p6);
  \draw[op] (p5.south) |- (p6.east);
  \draw[op] (p6) -- (d3);
  \draw[op] (d3) -- node[above, font=\scriptsize] {是} (t1);
  \draw[op] (d3) -- node[above, font=\scriptsize] {否} (t2);
\end{tikzpicture}
}
\end{center}

\subsection{正确性论证}

\textbf{为何"最大分量轴"投影能避免再次退化？}
设原三角形的法向 $\vec{n}$ 在某轴（如 $z$）上的分量最大，
则将三角形投影到 $xy$ 平面后，新法向 $\vec{n}' = (n_x, n_y)$
的长度 $\sqrt{n_x^2 + n_y^2}$ 是所有投影中\textbf{最大}的。
若 $\vec{n} \ne \vec{0}$（原三角形非退化），则至少一个分量非零，
该最大投影必非零；反之若 $\vec{n} = \vec{0}$（原三角形退化），
所有投影法向均为零，任意轴都判出 \texttt{orient2d $=$ 0}。

\textbf{为何不直接用 $|\vec{n}| < \varepsilon$？}
浮点阈值 $\varepsilon$ 难以选取：太小则接近退化的三角形漏报，
太大则正常窄三角形误报。Shewchuk 谓词给出\textbf{精确}符号，
无需阈值，能在边界情况下做出几何上正确的判定。

\section{与旧浮点实现的对比}

\begin{center}
\renewcommand{\arraystretch}{1.20}
\begin{tabular}{@{}lll@{}}
  \toprule
  \textbf{场景} & \textbf{旧实现（浮点阈值）} & \textbf{新实现（鲁棒谓词）} \\
  \midrule
  2D 三点方向      & \texttt{(b-a)$\times$(c-a) $>$ 0}              & \texttt{orient2d(a, b, c) $>$ 0} \\
  2D 共线判定      & $|cross| < 10^{-14}$                            & \texttt{orient2d $==$ 0} \\
  2D 三角形面积    & $\tfrac{1}{2}|cross|$                           & \texttt{signed\_area2d}（基于 \texttt{orient2d}） \\
  2D 点在三角形内  & 朴素叉积同号                                    & \texttt{point\_in\_triangle\_2d}（4 次 \texttt{orient2d}） \\
  3D 退化三角形    & $|cross| < 10^{-12}$                            & \texttt{is\_triangle\_degenerate\_3d} \\
  3D 四面体体积    & 朴素行列式                                      & \texttt{tet\_signed\_volume}（基于 \texttt{orient3d}） \\
  射线-三角形求交  & Möller-Trumbore（浮点行列式）                   & \texttt{orient3d} + 投影2D + \texttt{point\_in\_triangle\_2d} \\
  拉伸侧面退化     & $|edge \times offset| < 10^{-12}$               & \texttt{is\_triangle\_degenerate\_3d} \\
  网格自检退化面   & $area < 10^{-12}$                               & \texttt{is\_triangle\_degenerate\_3d} \\
  \bottomrule
\end{tabular}
\end{center}

\textbf{关键改进}：旧实现在退化边界（如点恰好共线、面积接近零）附近
依赖浮点阈值，存在误报 / 漏报风险；新实现符号判定\textbf{精确}，
仅在真正退化时报告，避免阈值选取的工程难题。

\section{调用方迁移清单}

本模块在 \texttt{halfedge} 库内的替换位置：

\begin{center}
\renewcommand{\arraystretch}{1.18}
\begin{tabular}{@{}lll@{}}
  \toprule
  \textbf{源文件} & \textbf{原函数} & \textbf{替换为} \\
  \midrule
  \texttt{triangulation.rs} & \texttt{v2\_cross / signed\_area}        & \texttt{orient2d / signed\_area2d} \\
  \texttt{triangulation.rs} & \texttt{point\_in\_triangle}             & \texttt{point\_in\_triangle\_2d} \\
  \texttt{triangulation.rs} & 凸顶点 \texttt{v2\_cross $\le 0$}        & \texttt{!is\_convex\_vertex2d} \\
  \texttt{geometry.rs}      & 体积手写行列式                            & \texttt{tet\_signed\_volume} \\
  \texttt{geometry.rs}      & Möller-Trumbore                          & \texttt{orient3d} + 投影2D \\
  \texttt{boolean.rs}       & 私有 \texttt{ray\_triangle\_intersection} & 调用 \texttt{geometry::ray\_triangle\_intersection} \\
  \texttt{validate.rs}      & \texttt{area < DEGENERATE\_AREA\_EPS}    & \texttt{is\_triangle\_degenerate\_3d} \\
  \texttt{topology\_ops.rs} & 拉伸侧面 \texttt{area $<$ 1e-12}（2 处） & \texttt{is\_triangle\_degenerate\_3d} \\
  \bottomrule
\end{tabular}
\end{center}

\section{退化守卫汇总}

\begin{center}
\renewcommand{\arraystretch}{1.18}
\begin{tabular}{@{}lll@{}}
  \toprule
  \textbf{API} & \textbf{退化场景} & \textbf{处理} \\
  \midrule
  \texttt{orient2d / orient3d / incircle / insphere}
    & 退化输入（共线 / 共面 / 共圆 / 共球）
    & 返回 \texttt{0.0}（精确零，非近似） \\
  \texttt{point\_in\_triangle\_2d}
    & 三角形退化（共线）       & 返回 \texttt{false} \\
  \texttt{point\_in\_triangle\_2d}
    & $p$ 在边上 / 顶点上       & 返回 \texttt{true}（含边界） \\
  \texttt{is\_triangle\_degenerate\_3d}
    & 三点重合                  & 返回 \texttt{true} \\
  \texttt{is\_triangle\_degenerate\_3d}
    & 法向某分量为零            & 自动选最大分量轴投影，避免再次退化 \\
  \texttt{tet\_signed\_volume}
    & 四点共面                  & 返回 \texttt{0.0} \\
  \bottomrule
\end{tabular}
\end{center}

\section{使用示例}

\begin{lstlisting}[language=Rust]
use halfedge::predicates::*;

// 1. 2D 方向判定
let a = [0.0, 0.0];
let b = [1.0, 0.0];
let c = [0.0, 1.0];
assert!(orient2d(a, b, c) > 0.0);  // CCW
assert!(is_ccw2d(a, b, c));
assert!((signed_area2d(a, b, c) - 0.5).abs() < 1e-15);

// 2. 浮点易错的共线判定（朴素 f64 会失败）
let p = [1e20, 1.0];
let q = [1e20 + 1.0, 1.0];
let r = [1e20 + 2.0, 1.0];
assert_eq!(orient2d(p, q, r), 0.0);  // Shewchuk 精确识别为共线

// 3. 2D 点在三角形内
assert!(point_in_triangle_2d([0.25, 0.25], a, b, c));
assert!(!point_in_triangle_2d([0.6, 0.6], a, b, c));

// 4. 3D 四面体体积
let a3 = [0.0, 0.0, 0.0];
let b3 = [1.0, 0.0, 0.0];
let c3 = [0.0, 1.0, 0.0];
let d3 = [0.0, 0.0, 1.0];
assert!((tet_signed_volume(a3, b3, c3, d3) - 1.0 / 6.0).abs() < 1e-15);

// 5. 3D 三角形退化判定
assert!(!is_triangle_degenerate_3d(a3, b3, c3));          // 正常三角形
assert!(is_triangle_degenerate_3d(a3, [1.0, 1.0, 1.0], [2.0, 2.0, 2.0]));  // 共线
assert!(is_triangle_degenerate_3d(a3, a3, a3));           // 重合
\end{lstlisting}

\section{测试覆盖}

\begin{center}
\renewcommand{\arraystretch}{1.18}
\begin{tabular}{@{}ll@{}}
  \toprule
  \textbf{测试名} & \textbf{验证点} \\
  \midrule
  \texttt{orient2d\_ccw}                       & CCW 三点 \texttt{orient2d $>$ 0}，\texttt{is\_ccw2d} 为真 \\
  \texttt{orient2d\_cw}                        & CW 三点 \texttt{orient2d $<$ 0} \\
  \texttt{orient2d\_collinear\_exact}          & 精确共线三点 \texttt{orient2d $=$ 0} \\
  \texttt{orient2d\_collinear\_near}           & 浮点易错共线（$10^{20}$ 量级），仍精确返回 0 \\
  \texttt{signed\_area\_2d\_unit\_triangle}    & 单位三角形面积 $= 0.5$ \\
  \texttt{point\_in\_triangle\_interior}       & 内部点返回 \texttt{true} \\
  \texttt{point\_in\_triangle\_vertex}         & 三顶点返回 \texttt{true} \\
  \texttt{point\_in\_triangle\_edge}           & 边上点返回 \texttt{true} \\
  \texttt{point\_outside\_triangle}            & 外部点返回 \texttt{false} \\
  \texttt{point\_in\_triangle\_cw\_orientation}& CW 三角形仍能正确判定内点 \\
  \texttt{point\_in\_triangle\_degenerate}     & 共线三角形返回 \texttt{false} \\
  \texttt{orient3d\_basic}                     & Shewchuk 约定：下方 $>$ 0 / 上方 $<$ 0 / 共面 $=$ 0 \\
  \texttt{tet\_signed\_volume\_unit}           & 标准正交四面体体积 $= 1/6$ \\
  \texttt{orient3d\_coplanar}                  & 四点共面 \texttt{is\_coplanar} 为真 \\
  \texttt{incircle\_inside / outside / on\_circle} & 共圆谓词三类情况 \\
  \texttt{insphere\_inside / outside}          & 共球谓词（要求 abcd 正方向） \\
  \texttt{convex\_vertex\_ccw}                 & CCW 凸顶点 / CW 凹顶点 \\
  \texttt{triangle\_degenerate\_3d\_normal}    & 正常 3D 三角形非退化 \\
  \texttt{triangle\_degenerate\_3d\_collinear} & 3D 共线三点判为退化 \\
  \texttt{triangle\_degenerate\_3d\_coincident}& 三点重合判为退化 \\
  \texttt{triangle\_degenerate\_3d\_axis\_aligned} & 主轴对齐三角形的退化 / 非退化 \\
  \bottomrule
\end{tabular}
\end{center}

\texttt{src/predicates.rs} 共 \textbf{20 个}单元测试，全库共 \textbf{433 个}。

\section{参考资料}

\begin{itemize}
  \item Shewchuk, J. R. (1997). \textit{Adaptive Precision Floating-Point
        Arithmetic and Fast Robust Geometric Predicates.}
        Discrete \& Computational Geometry, 18(3), 305--363.
  \item \texttt{robust} crate 文档：\url{https://docs.rs/robust}
  \item Möller-Trumbore 射线求交原始论文（已被本模块替换）：
        Möller, T., \& Trumbore, B. (1997). Fast, Minimum Storage
        Ray-Triangle Intersection. \textit{Journal of Graphics Tools}, 2(1).
\end{itemize}

\end{document}
